\section{总结与延伸}

本节把整期压缩成六个结论。第一，机器人革命的关键是从专用模型走向通用模型。第二，早期路线是用 LLM/VLM/Code LM 改造 Planning、Perception、Actuation。第三，真正的 VLA 要把 Vision、Language、Action 放进端到端模型。第四，RT-2/OpenVLA 等路线借助 VLM 迁移互联网知识，但 action 高频控制仍需要专门处理。第五，Diffusion Policy/RDT 说明动作本身可以用生成模型建模。第六，未来方向是世界模型、统一理解预测和在线强化学习。

\begin{importantbox}{把 VLA 投屏课放进张小珺 AI/互联网队列}
这期是 EP106、EP109、EP121 等具身智能访谈的技术底座。它不是创业故事，而是机器人基座模型路线图，解释了为什么 VLA 会成为 2025 年机器人讨论里的关键词。
\end{importantbox}

\subsection{经典论文速查表}

\noindent\begin{tabularx}{\textwidth}{p{0.20\textwidth}p{0.30\textwidth}X}
\toprule
论文/系统 & 关键贡献 & 本课位置 \\
\midrule
SayCan & LLM 规划 + affordance 可行性 & 用语言模型替代 Planning 的代表。 \\
Inner Monologue & 环境反馈和语言化自我修正 & 让执行反馈回到规划。 \\
DoReMi & 检测和恢复计划执行错配 & 处理真实机器人执行失败。 \\
VoxPoser & 3D value maps & 把语言目标转成空间操作场。 \\
ALOHA & 低成本双臂操作和 ACT & 示范数据和动作序列学习。 \\
RT-1/RT-2 & 机器人 Transformer / VLA 迁移 & Google Robotics 主线。 \\
RT-X/OpenVLA & 跨本体数据和开源 VLA & 数据与开源基础设施。 \\
HiRT/pi0/GR00T N1 & VLM + Action Head/Expert & 更认真处理 action。 \\
Diffusion Policy/RDT & 动作扩散和双手操作基础模型 & 生成式 action policy。 \\
UP-VLA/iRe-VLA & 统一理解预测和在线 RL & 未来方向。 \\
\bottomrule
\end{tabularx}

\subsection{后续观察问题}

最后回到张小珺这组 AI/互联网访谈的长期观察框架：VLA 不是一个单点模型名，而是一组围绕数据、模型、硬件、商业落地和安全训练的系统问题。下面这些问题适合作为后续追踪清单，用来判断机器人基座模型是否真的进入类似大语言模型的 scale 阶段。

\begin{enumerate}
\item VLA 的 scaling law 是否会像语言模型一样稳定出现？
\item 跨本体数据集能否足够大、足够标准化，支撑真正通用策略？
\item VLM backbone 直接输出 action，和 VLM + Action Head，哪条路线更可扩展？
\item 世界模型是否会成为机器人泛化的关键，还是 action policy 本身就足够？
\item 在线强化学习能否在真实机器人上安全、高效、可持续地运行？
\item 人形机器人是否会先靠本体和遥操作落地，还是必须等待 VLA 能力成熟？
\end{enumerate}

\subsection{拓展阅读}

\begin{itemize}
\item 对具身智能学术史感兴趣，可对照 EP106 王鹤访谈。
\item 对仿真和合成数据感兴趣，可对照 EP109 谢晨/光轮访谈。
\item 对多模态世界模型和视觉推理感兴趣，可对照 EP102 张祥雨访谈。
\item 对机器人产业化和硬件入口感兴趣，可对照 EP104 Rokid 与 EP121 谭捷访谈。
\end{itemize}

\end{document}
